% Part: conditional-logics
% Chapter: minimal-change-semantics
% Section: true-false

\documentclass[../../../include/open-logic-section]{subfiles}

\begin{document}

\olfileid[pt]{con}{min}{tf}

\olsection{Verdade e falsidade dos contrafactuais}

Um contrafactual $!A \cif !B$ é verdadeiro de modo não vacuoso se os
mundos-$!A$ mais próximos são todos mundos-$!B$, como mostra \olref{fig:true}.
\begin{figure}
\begin{center}
\begin{tikzpicture}[scale=.7]\small
  \spheresystem{5}
  \draw (0,0) node {$w$};
  \propositionintersect{0}{4}{40}{3.5}
  {\tikzset{proposition/.append={smooth,tension=1.4}}
  \proposition[shift={(1.6,-1)}]{90}{1}{30}{4}}
  \path (1.5,3) node[below] {$\formula{B}$};
  \path (3,0) node {$\formula{A}$};
\end{tikzpicture}
\caption{Contrafactual verdadeiro de modo não vacuoso}
\ollabel{fig:true}
\end{center}
\end{figure}
Um contrafactual também é verdadeiro em $w$ se o sistema de esferas
ao redor de~$w$ não contém nenhuma esfera que admita $!A$. Nesse caso, ele é
verdadeiro \emph{vacuamente} (veja \olref{fig:vacuous}).
\begin{figure}
\begin{center}
\begin{tikzpicture}[scale=.7]\small
  \spheresystem{5}
  \draw (0,0) node {$w$};
  \proposition{0}{6.5}{40}{3.5}
  {\tikzset{proposition/.append={smooth,tension=1.4}}
  \proposition[shift={(1.2,-1)}]{90}{1}{30}{4}}
  \path (1.5,3) node[below] {$\formula{B}$};
  \spherepos{0}{7}{node {$\formula{A}$}}
\end{tikzpicture}
\caption{Contrafactual verdadeiro vacuamente}
\ollabel{fig:vacuous}
\end{center}
\end{figure}

Ele pode ser falso de duas maneiras. Uma delas ocorre quando os mundos-$!A$
mais próximos não são todos mundos-$!B$, embora alguns o sejam. Nesse caso, $!A \cif
\lnot !B$ também é falso (veja \olref{fig:false}).
\begin{figure}
\begin{center}
\begin{tikzpicture}[scale=.7]\small
  \spheresystem{5}
  \draw (0,0) node {$w$};
  \propositionintersect{0}{4}{40}{3.5}
  {\tikzset{proposition/.append={smooth,tension=1.4}}
  \proposition[shift={(1.6,0)}]{90}{1}{30}{3}}
  \path (1.5,3) node[below] {$\formula{B}$};
  \path (3,0) node {$\formula{A}$};
\end{tikzpicture}
\caption{Contrafactual falso, oposto falso}
\ollabel{fig:false}
\end{center}
\end{figure}
Se os mundos-$!A$ mais próximos não se sobrepõem de modo algum aos mundos-$!B$,
então $!A \cif !B$ é falso. Mas, nesse caso, todos os mundos-$!A$ mais próximos são
mundos-$\lnot !B$, e portanto $!A \cif \lnot !B$ é verdadeiro (veja
\olref{fig:false-opposite}).
\begin{figure}
\begin{center}
\begin{tikzpicture}[scale=.7]\small
  \spheresystem{5}
  \draw (0,0) node {$w$};
  \propositionintersect{0}{4}{40}{3.5}
  {\tikzset{proposition/.append={smooth,tension=1.4}}
  \proposition[shift={(-1,0)}]{90}{1}{30}{3}}
  \path (-1,3) node[below] {$\formula{B}$};
  \path (3,0) node {$\formula{A}$};
\end{tikzpicture}
\caption{Contrafactual falso, oposto verdadeiro}
\ollabel{fig:false-opposite}
\end{center}
\end{figure}

Em contraste com o condicional estrito, os contrafactuais podem ser
contingentes. Considere o modelo de esferas em \olref{fig:contingent}. Os
mundos-$!A$ mais próximos de~$u$ são todos mundos-$!B$, portanto $\mSat{M}{!A \cif
  !B}[u]$. Mas há mundos-$!A$ mais próximos de~$v$ que não são
mundos-$!B$, portanto $\mSat/{M}{!A \cif !B}[v]$.
\begin{figure}
\begin{center}
\begin{tikzpicture}[scale=.6]\small
  \spheresystem{7}
  \spheresystem[shift={(0:3.1)},dashed]{4}
  \propositionintersect{45}{4}{40}{4.5}
  \begin{scope}
    \clip \propositionplot{45}{4}{40}{4.5} ;
    \spherelayer[shift={(0:3.1)}]{3}
  \end{scope}
  \proposition[shift={(1.4,-.5)}]{90}{1}{35}{5}
  \path (0,0) node {$u$};
  \path (0:3.1) node {$v$};
  \spherepos{35}{9}{node {$\formula{A}$}}
  \spherepos{80}{9}{node {$\formula{B}$}}
\end{tikzpicture}
\end{center}
\caption{Contrafactual contingente}
\ollabel{fig:contingent}
\end{figure}
\end{document}
